%% conclusion.tex
%%

%% ==================
\section{Conclusion and Future Work}
\label{sec:Conclusion}
%% ==================

This project asked whether a federated learning alternative to Et-Tousy et al.'s (2026) centralized Random Forest -- one model per IoT site trained on local data only, combined without pooling raw telemetry -- could recover most of the centralized model's accuracy while resolving the privacy/scalability limitation the paper names but does not build. Across 5 seeds on synthetic, non-IID multi-site QoS telemetry, the federated ensemble reached 99.49\% macro-F1 versus the centralized baseline's 99.96\% and a single site's own local-only model's 97.37\% -- recovering the large majority of the benefit of centralization without requiring it.

The clearest finding is not the small centralized-vs-federated gap, but the large local-only-vs-federated gap: a site that cannot collaborate with any other site (no data sharing, no federation) generalises measurably worse to the traffic patterns it will actually see in production. Federation is the difference between a site being stuck with a locally-biased model and getting most of the benefit of the full population's data, without that data ever leaving its site of origin.

\textbf{Future work:}
\begin{itemize}
  \item Evaluate on real multi-site OneM2M/Azure IoT telemetry rather than synthetic data, to check whether the federation benefit holds under real sensor noise and real non-IID splits.
  \item Extend the federated design from simple probability-averaging to a weighted aggregation scheme (e.g. weighting each site's vote by its local validation confidence), which the current implementation does not attempt.
  \item Evaluate secure aggregation (so that not even model updates are individually inspectable by a central coordinator), extending the privacy property this project targets but does not fully secure.
\end{itemize}
